\begin{table}[H]
\centering
\caption{MFCC e LDA na placa, com os segmentos já decimados a 12,8~kHz mandados pela serial (a decimação fica fora desta rodada; na placa, ela é testada pelo recorte de 1~s a 51,2~kHz do segmento 29, no boot da mesma captura, e no PC, nas 5 gravações inteiras): concordância com o Python e diferença do escore em fração da tolerância encadeada ($10^{-5}$ relativo, contra o float64). No Jung, todos os segmentos são do treino do modelo final: mede a fidelidade do porte, e não o desempenho.}
\label{tab:serial_placa_resumo}
\small
\begin{tabular}{lrrrrl}
\toprule
\textbf{Conjunto} & \textbf{Processados} & \textbf{Concordância} & \textbf{\% da tol. (mediana)} & \textbf{\% da tol. (máx.)} & \textbf{Pior segmento} \\
\midrule
Jung & 295/295 & 295/295 & 2,2 & 9,8 & \texttt{bpfi\_1.0mm\_seg0127} \\
HUSTmotor & 40/40 & 40/40 & 1,3 & 7,6 & \texttt{bf\_20hz\_a4} \\
\bottomrule
\end{tabular}
\end{table}

\begin{table}[H]
\centering
\caption{Matriz de confusão do sistema embarcado (segmentos pela serial) e das decisões do Python. No HUSTmotor, a LDA do Jung sem retreino (\texttt{exp274} do Protocolo C).}
\label{tab:serial_placa_matriz}
\small
\begin{tabular}{llrrrr}
\toprule
 & & \multicolumn{2}{c}{\textbf{Placa}} & \multicolumn{2}{c}{\textbf{Python}} \\
\cmidrule(lr){3-4}\cmidrule(lr){5-6}
\textbf{Conjunto} & \textbf{Classe verdadeira} & normal & falha & normal & falha \\
\midrule
Jung & normal & 59 & 0 & 59 & 0 \\
 & falha & 0 & 236 & 0 & 236 \\
\midrule
HUSTmotor & normal & 0 & 20 & 0 & 20 \\
 & falha & 0 & 20 & 0 & 20 \\
\bottomrule
\end{tabular}
\end{table}
